\begin{frame}{¿Qué es Git?}
  \begin{columns}[onlytextwidth]
    \column{0.5\textwidth}
    \begin{center}
      \includegraphics[scale=0.07]{images/logoGit2}
    \end{center}
    \column{0.5\textwidth}
    Git es un Sistema de Control de Versiones (VCS) que nos permite gestionar los cambios que realizamos un proyecto.
  \end{columns}
\end{frame}

\begin{frame}{¿Qué es Git?}
  \only<1>{\begin{center}
    \includegraphics[scale=0.9]{images/caos}
  \end{center}}
  \only<2>{\begin{center}
    \includegraphics[scale=0.9]{images/nocaos}
  \end{center}}
\end{frame}

\begin{frame}{¿Git == GitHub?}
  \begin{columns}[onlytextwidth]
    \column{0.25\textwidth}
    \begin{center}
      \animategraphics[height=1.3in,autoplay,loop] {12}{images/DF/githubDP1_}{0}{7}
    \end{center}
    \column{0.5\textwidth}
    \begin{center}
      \alert{\huge ¿Git == GitHub? \\
      \smallskip
      \uncover<2->{
      {\color{solarized-red} \textbf{NO}}}} \\
      \bigskip
      \uncover<3>{
      GitHub es una forja (plataforma de desarrollo colaborativo) para alojar proyectos Git.}
    \end{center}
    \column{0.25\textwidth}
    \begin{center}
      \animategraphics[height=1.3in,autoplay,loop] {12}{images/DF/githubDP2_}{0}{13}
    \end{center}
  \end{columns}
\end{frame}

\begin{frame}{Otras forjas}
  \begin{columns}[onlytextwidth]
    \column{0.5\textwidth}
    \begin{center}
      \includegraphics[scale=0.3]{images/logoGitLab}
    \end{center}
    \column{0.5\textwidth}
  \end{columns}
  \begin{columns}[onlytextwidth]
    \column{0.5\textwidth}
    \column{0.5\textwidth}
    \begin{center}
      \includegraphics[scale=0.4]{images/logoBitbucket}
    \end{center}
  \end{columns}
  \begin{columns}[onlytextwidth]
    \column{0.5\textwidth}
    \begin{center}
      \includegraphics[scale=0.25]{images/logoSourceForge}
    \end{center}
    \column{0.5\textwidth}
  \end{columns}
\end{frame}

\begin{frame}{¿Cómo se utiliza Git?}
  % \alert{\Large ¿Cómo se utiliza Git?}
  \begin{columns}[onlytextwidth]
    \column{0.6\textwidth}
    Existen:
    \begin{itemize}
      \item Clientes gráficos (GitKraken, SourceTree, GitHub Desktop...)
      \item Integración en muchos editores e IDEs (Atom, Visual Studio Code, IntelliJ IDEA...)
    \end{itemize}
    \vspace{0.5cm}
    \uncover<2>{
    {\color{violet}\large Git se utiliza mediante comandos en la \textbf{Terminal}.}}
    \column{0.4\textwidth}
    \uncover<2>{
    \includegraphics[scale=0.4]{images/hackerMan.eps}}
  \end{columns}
  \uncover<1>{
    \begin{tikzpicture}[remember picture,overlay]
      \node[anchor=north east, inner xsep=30pt, inner ysep=40pt] at (current page.north east) {\includegraphics[scale=0.11]{images/logoGitKraken.png}};
      \node[anchor=north east, inner xsep=80pt, inner ysep=50pt] at (current page.north east) {\includegraphics[scale=0.3]{images/logoSourceTree.png}};
      \node[anchor=north east, inner xsep=30pt, inner ysep=120pt] at (current page.north east) {\includegraphics[scale=0.03]{images/logoAtom.png}};
      \node[anchor=north east, inner xsep=80pt, inner ysep=180pt] at (current page.north east) {\includegraphics[scale=0.045]{images/logoVisualStudioCode.png}};
      \node[anchor=north east, inner xsep=130pt, inner ysep=135pt] at (current page.north east) {\includegraphics[scale=0.040]{images/logoIntelliJIDEA.png}};
    \end{tikzpicture}
  }
\end{frame}

\begin{frame}{¿Cómo funciona Git? - Las 3 áreas}
  \begin{center}
    \begin{tikzpicture}[head/.style={line width=1pt, draw=solarized-base01, rounded corners=1pt, text width=3cm, align=center}]
    \node[head, fill=solarized-orange!25] (wd) {Working directory};
    \node[head, fill=solarized-blue!25, right=of wd] (sa) {Staging \\ area};
    \node[head, fill=solarized-base1!25, right=of sa] (gr) {.git directory \\ (Repository)};
    \begin{scope}[line width=2pt,solarized-base01]
    \foreach \box in {wd,sa,gr}
      \draw (\box.south) -- +(0,-4cm);
    \end{scope}
    \begin{scope}[every node/.style={single arrow, line width=1pt, draw=solarized-base01, fill=solarized-green!25, anchor=west, align=center}]
    \node [anchor=east, text width=7.8cm,shape border rotate=180] at ([yshift=-1cm]gr.south) {Checkout the project};
    \node [anchor=west, text width=3.5cm] at ([yshift=-2.2cm]wd.south) {Staged fixed};
    \node [anchor=west, text width=3.5cm] at ([yshift=-3.4cm]sa.south) {Commit};
    \end{scope}
    \end{tikzpicture}
  \end{center}
\end{frame}

\begin{frame}{¿Cómo funciona Git? - Estados de los archivos}
  \begin{center}
    \begin{tikzpicture}[head/.style={line width=1pt, draw=solarized-base01, rounded corners=1pt, text width=2.3cm, align=center}]
    \node[head, fill=solarized-base1!25] (ut) {Untracked};
    \node[head, fill=solarized-blue!25, right=of ut] (um) {Unmodified};
    \node[head, fill=solarized-yellow!25, right=of um] (mo) {Modified};
    \node[head, fill=solarized-orange!25, right=of mo] (st) {Staged};
    \begin{scope}[line width=2pt,solarized-base01]
    \foreach \box in {ut,um,mo,st}
      \draw (\box.south) -- +(0,-6cm);
    \end{scope}
    \begin{scope}[every node/.style={single arrow, line width=1pt, draw=solarized-base01, fill=solarized-green!25, anchor=west, align=center}]
    \node [anchor=west, text width=10.01cm] at ([yshift=-1cm]ut.south) {Add the file};
    \node [anchor=west, text width=2.8cm] at ([yshift=-2.1cm]um.south) {Edit the file};
    \node [anchor=west, text width=2.8cm] at ([yshift=-3.2cm]mo.south) {Stage the file};
    \node [anchor=east, text width=2.8cm,shape border rotate=180] at ([yshift=-4.3cm]um.south) {Remove the file};
    \node [anchor=east, text width=6.4cm,shape border rotate=180] at ([yshift=-5.4cm]st.south) {Commit};
    \end{scope}
    \end{tikzpicture}
  \end{center}
\end{frame}
